% Appendix C — 9 CLEAR-grounded prompt templates (compact, verbatim).
% Reproducibility: identical templates live in p1/python/src/prompt_builder.py.

\section*{Appendix C: Prompt Templates}\label{sec:appendix-c}

Nine templates (three CLEAR levels~$\times$ three pedagogical framings) were
applied verbatim; \texttt{\{signature\}} is replaced by the HumanEval
\texttt{prompt} field (function signature plus docstring). Source of truth:
\texttt{p1/python/src/prompt\_builder.py}.

{\footnotesize
\begin{description}
  \item[Algorithm explanation --- L1 basic.] \texttt{Explain this:
  \{signature\}}

  \item[Algorithm explanation --- L2 intermediate.] You are teaching CS.
  Concisely explain the algorithm for this function. List the steps in
  order.~\texttt{\{signature\}}

  \item[Algorithm explanation --- L3 advanced.] As a CS tutor for an IT
  undergraduate, provide a structured explanation of the algorithm behind
  the function below. Follow these steps: (1) restate the problem in plain
  English; (2) outline the high-level approach; (3) walk through the key
  invariants or data flow; (4) state when this approach is preferred over
  alternatives. Use concise prose, adapt vocabulary to an undergraduate
  audience, and end with a one-sentence self-check:
  `Does this explanation help someone solve a similar problem
  themselves?'~\texttt{\{signature\}}

  \item[Complexity analysis --- L1 basic.] \texttt{What is the complexity
  of this? \{signature\}}

  \item[Complexity analysis --- L2 intermediate.] Analyze the time and
  space complexity of this function. State the Big-O at the
  end.~\texttt{\{signature\}}

  \item[Complexity analysis --- L3 advanced.] As a CS tutor, conduct a
  complexity analysis for the function below. Follow these steps: (1)
  identify loops and recursions; (2) count operations symbolically;
  (3) derive worst-case time $T(n)$ and auxiliary space $S(n)$;
  (4) comment on whether this is asymptotically optimal for the problem;
  (5) end with a self-check: `Can a student replicate this analysis on a
  similar problem?'~\texttt{\{signature\}}

  \item[Problem solving --- L1 basic.] \texttt{Solve: \{signature\}}

  \item[Problem solving --- L2 intermediate.] Write a Python
  implementation for the function below. Return only the function body
  (no explanation).~\texttt{\{signature\}}

  \item[Problem solving --- L3 advanced.] You are a senior Python engineer
  mentoring an IT undergraduate. Write a clean, idiomatic Python
  implementation for the function below. Requirements: (1) include a
  one-line docstring summary; (2) use Pythonic constructs where
  appropriate; (3) avoid external libraries; (4) aim for the best
  complexity you can justify; (5) end with a brief comment verifying edge
  cases you considered.~\texttt{\{signature\}}
\end{description}
}
